\chapter*{Motivație} 
\addcontentsline{toc}{chapter}{Motivație}

Pe parcursul dezvoltării discipline informaticii, natura a servit ca inspirație pentru oamenii de știința care se ocupau cu dezvoltarea de noi algoritmi. Astfel a aparut clasa de 
\textbf{algoritmi evolutivi}, din care se branșează mai departe alte familii de algoritmi:
\textbf{algoritmii genetici}, \textbf{strategii evolutive}, \textbf{programarea genetica} și \textbf{algoritmii de neuro-evoluție}.\\
Aceasta familie de algoritmi a ajutat la dezvoltarea multor solutii unde algoritmii clasici deterministi sunt prea costisitori sau ineficienti. Un exemplu ar fi o antena X-band evoluata pentru misiunea NASA din 2006 numita Space Technology 5, al carei design a fost gasit prin intermediul unui algoritm evolutiv\cite{Space2006Antenna}.\\
Pe parcursul timpului, studiul in natura a intercatiunilor dintre prada si pradator a permis extragerea de informatii valoroase despre dinamica ecosistemelor.
Datorita avansului rapid al inteligentii artificiale si a algoritmilor evolutivi, avem oportunitatea de a cerceta aceste intercatiuni intr-un mediu virtual.\\
Aceasta teza de licenta propune utilizare algoritmilor \textbf{NEAT}(NeuroEvolution of Augmenting Topologies) si \textbf{rtNEAT}(real time NEAT), plus mediul de dezvoltare de jocuri \textbf{Unity}(2D) pentru crearea unei simulari dinamice a intercatiunilor dintre speciile de prada si cele pradator. Intelegerea acestor intercatiuni este esentiala pentru a descoperi factorii care influențează dinamica populațiilor, coexistența speciilor și stabilitatea ecosistemelor. Utilizand un mediu virtual, putem studia interactiunile prada-pradator in izolatie, putem controla variabilele de mediu, si putem repeta experimente de mai multe ori intr-un mod consistent(relativ la variabilitatea naturii).
Algoritmul NEAT, datorita capacitatii sale de a evolua atat ponderile cat si topologiile rețelelor neuronale, reprezinta un  motor puternic pentru modelarea comportamentelor și proceselor de luare a deciziilor ale prădătorilor și prăzii, cu scopul de a supravetui mai mult. Prin evoluția rețelelor neurale ale prăzii, putem explora comportamentele de căutare a hranei, tactici de evitare și strategii de supraviețuire, iar pradatorii ar trebui sa dezvolte strategii de vanatoare, cu adaptari pentru schimbarile din mediu.\\
Motorul de joc Unity permite sa vizualizam in timp real comportamentele pe care le dezvolta orgnismele, sa monitorizam noile specii care apar din ambele categorii de organism, sa adaugam potential noi obstacole, noi senzori, sau noi abilitati prazii si pradatorilor, si sa controlam variabilele si viteza simularii.\\
În concluzie, această teză de licență își propune să îmbine domeniile inteligenței artificiale, al evolutiei biologice si al ecologiei, utilizând algoritmul NEAT și Unity pentru a simula interacțiunile dintre pradă și prădători. Prin studierea acestor dinamici ne propunem sa contribuim intelegerea sistemelor ecologice, cat și la educare populatiei pe larg asupra acestor interactiuni prin intermediul vizualizarilor. Ipoteza noastra e ca studiind evolutia populatiilor prazii si pradatorilo aceste se vor conforma ecuatiilor Lotka-Voltera.

